\documentclass[UTF8,zihao=-4]{ctexart}
\usepackage[a4paper,margin=2.2cm]{geometry}
\usepackage{xcolor}
\usepackage{hyperref}
\usepackage{enumitem}
\usepackage{fontspec}
\usepackage{amsmath}
\IfFontExistsTF{Microsoft YaHei}{\setCJKmainfont{Microsoft YaHei}}{\setCJKmainfont{SimSun}[BoldFont=SimHei]}
\IfFontExistsTF{Microsoft YaHei UI}{\setCJKsansfont{Microsoft YaHei UI}}{}
\IfFontExistsTF{Consolas}{\setmonofont{Consolas}}{}
\definecolor{linkblue}{RGB}{30,80,145}
\hypersetup{colorlinks=true,linkcolor=linkblue,urlcolor=linkblue,pdfborder={0 0 0}}
\setlength{\parindent}{2em}
\setlength{\parskip}{0.25em}
\linespread{1.16}
\setlist{itemsep=0.2em,topsep=0.35em,leftmargin=2.5em}
\pagestyle{plain}

\title{06\_快排平均复杂度递推分析}

\author{}

\date{}

\begin{document}

\maketitle

\section*{06 快排平均复杂度递推分析}

\subsection*{题目原文}

复用快排中平均复杂分析的方法来处理其它问题：

\[A(n) = (n - 1) + (2/n) \sum_{i=1}^{n-1} A(i)\]

\subsection*{考查类型}

递推分析题。

\subsection*{递推来源}

快速排序平均复杂度分析中，划分一次需要 \texttt{n - 1} 次比较。若主元位置在所有位置上等概率出现，就会得到左右子问题规模的平均递推。

PDF 中给出的递推是：

\[A(n) = (n - 1) + (2/n) \sum_{i=1}^{n-1} A(i)\]

通常取：

\[A(1) = 0\]

\subsection*{解法：乘 n 再相减}

原式两边乘以 $n$：

\[nA(n) = n(n - 1) + 2 \sum_{i=1}^{n-1} A(i)\]

对 \texttt{n - 1} 写同样的式子：

\[(n - 1)A(n - 1) = (n - 1)(n - 2) + 2 \sum_{i=1}^{n-2} A(i)\]

两式相减：

\begin{quote}
\small\ttfamily\raggedright
\relax nA(n)\ -\ (n\ -\ 1)A(n\ -\ 1)\\
\end{quote}
\[= 2(n - 1) + 2A(n - 1)\]

整理得：

\[nA(n) = (n + 1)A(n - 1) + 2(n - 1)\]

两边除以 \texttt{n(n + 1)}：

\[A(n)/(n + 1) = A(n - 1)/n + 2(n - 1)/(n(n + 1))\]

令：

\[B(n) = A(n)/(n + 1)\]

则：

\[B(n) = B(n - 1) + 2(n - 1)/(n(n + 1))\]

累加得到：

\[A(n) = 2(n + 1)H_{n} - 4n\]

其中：

\[H_{n} = 1 + 1/2 + ... + 1/n\]

由于：

\[H_{n} = \Theta(\log n)\]

所以：

\[A(n) = \Theta(n \log n)\]

\subsection*{答题模板}

\begin{enumerate}
\item 写出平均递推。
\item 两边乘以 $n$，消去分母。
\item 写出 \texttt{n - 1} 的同类递推。
\item 两式相减，去掉求和。
\item 变形为可累加形式。
\item 用调和级数 $H_n = \Theta(\log n)$ 得出结果。
\end{enumerate}

\subsection*{例题}

\subsubsection*{例题 1}

题目：设：

\[A(1) = 0\]
\[A(n) = (n - 1) + (2/n) \sum_{i=1}^{n-1} A(i)\]

计算 \texttt{A(2)}、\texttt{A(3)}、\texttt{A(4)}。

解答：

\[A(2) = 1 + (2/2)A(1) = 1\]

\[A(3) = 2 + (2/3)(A(1) + A(2))\]
\[= 2 + (2/3)(0 + 1)\]
\[= 8/3\]

\[A(4) = 3 + (2/4)(A(1) + A(2) + A(3))\]
\[= 3 + (1/2)(0 + 1 + 8/3)\]
\[= 29/6\]

\subsubsection*{例题 2}

题目：用闭式公式验证 \texttt{A(4)}。

解答：

已知：

\[A(n) = 2(n + 1)H_{n} - 4n\]

其中：

\[H_{4} = 1 + 1/2 + 1/3 + 1/4 = 25/12\]

代入：

\[A(4) = 2\cdot 5\cdot (25/12) - 16\]
\[= 250/12 - 192/12\]
\[= 29/6\]

与递推计算结果一致。

\subsection*{常见误区}

不要直接把递推看成主定理形式。这里有求和项，不是 $T(n) = aT(n/b) + f(n)$。

不要漏掉 $H_n$ 的渐进阶。结论的关键是调和级数带来的 $\log n$。

\end{document}
